\appendix

\section{Finite-size gap analysis}
\label{sec:gap}

We take the gap in the $P=-2$ symmetry sector,
\begin{equation}
\Delta_{\rm sec}(L,s)=E_1^{(P=-2)}(L,s)-E_0^{(P=-2)}(L,s),
\end{equation}
where $E_0^{(P=-2)}$ ($E_1^{(P=-2)}$) is the ground-state (first-excited-state)
energy in that sector.
The sector gap at $\delta=0$ for $L=8,12,16$ is shown in
Fig.~\ref{fig:D02a}. Near criticality finite-size scaling gives
$\Delta_{\rm sec}L\approx A$; plotting $A$ versus $s$
(Fig.~\ref{fig:D02b}) the crossings for increasing $L$ stabilize at
$s=0.5$. Specifically, the drift of the crossing between successive sizes
shrinks as $L$ grows, consistent with convergence to $s=0.5$ in the
thermodynamic limit. The gap at $s=0.5$ versus $1/L$
($L=8,12,16,20,24\ $) is shown in Fig.~\ref{fig:D02c}; the linear fit through $L=20,24$ (coefficients
marked in the figure) has a small fitted intercept $b=0.05$, consistent with
the gap-closing trend. Finite-size effects therefore shift only the
phase boundaries slightly and leave the bulk phase structure
unaffected.
Ground states throughout are obtained by exact diagonalization
(Lanczos), and the crossing analysis follows standard finite-size
practice for quantum spin systems, for which we follow the lecture
notes of Ref.~\cite{sandvik2010computational}.

\begin{figure}[htbp]
\centering
\includegraphics[width=\columnwidth]{exp03_D02a.pdf}
\caption{Symmetry-sector gap $\Delta_{\rm sec}$ at $\delta=0$ for
$L=8,12,16$. \label{fig:D02a}}
\end{figure}

\begin{figure}[htbp]
\centering
\includegraphics[width=\columnwidth]{exp03_D02b.pdf}
\caption{$A=\Delta_{\rm sec}L$ versus $s$; crossings stabilize at
$s=0.5$ as $L\to\infty$. \label{fig:D02b}}
\end{figure}

\begin{figure}[htbp]
\centering
\includegraphics[width=\columnwidth]{exp03_D02c.pdf}
\caption{Gap at $s=0.5$ versus $1/L$ with linear fit through
$L=20,24$ (small fitted intercept $b=0.05$). \label{fig:D02c}}
\end{figure}

% CIT-003 (optional): finite-size scaling direction (bind)
